\chapter{Introducción}
\thispagestyle{fancy}
\label{ch:introducción}

Las características físicas de los materiales, como la densidad, son
 fundamentales para comprender su comportamiento y aplicaciones. 
 La densidad es una propiedad intensiva que relaciona la masa de
  un material con su volumen, y es crucial en diversas áreas de la 
  ciencia y la ingeniería. Al ser una propiedad intrínseca, la densidad no depende de la cantidad de
    material presente, lo que la convierte en una herramienta valiosa para identificar sustancias 
    y predecir su comportamiento en diferentes condiciones.
  
  En esta práctica, se lleva a cabo la 
  medición de la densidad de sólidos utilizando métodos experimentales híbridos;
  se obtiene la masa y dimensiones de los objetos de manera directa usando 
  instrumentos de medición y luego se calcula el volumen para determinar 
  la densidad. Este cálculo geométrico permite obtener la densidad sin necesidad de 
  sumergir los objetos en un líquido, 
  lo que es especialmente útil para materiales que pueden reaccionar entre ellos.

  Los instrumentos de medición utilizados durante la experiencia incluyen una balanza,
  un vernier y un micrómetro. La balanza se emplea para medir la masa de los
objetos con precisión, mientras que el vernier y el micrómetro se utilizan
 para medir las dimensiones de los objetos, 
lo que permite calcular su volumen.

Se espera que los resultados obtenidos en esta 
práctica sean consistentes con los valores teóricos 
de densidad para los materiales estudiados, lo que permitirá 
validar la precisión de los métodos experimentales utilizados. 
Además, se analizarán las posibles fuentes de error y se discutirán 
las implicaciones de los resultados obtenidos.
